% ============================================================================
% CONCLUSION
% ============================================================================

\section{Conclusion}
\label{sec:conclusion}

We present \ours{}, a production-ready framework that addresses the fundamental scalability limitations of existing LLM routing systems. By introducing shippable priors---compact covariance matrices distilled from offline expert supervision---we resolve the cold-start inefficiency that has historically prevented contextual bandit adoption in cost-sensitive deployments. Our approach achieves a \textbf{64.6\% reduction in day-one regret} while maintaining the adaptive plasticity necessary for evolving model landscapes.

\paragraph{Contributions to Production ML Systems.}
\ours{} makes three architectural contributions relevant to production machine learning systems beyond LLM routing. First, our shippable prior mechanism demonstrates a general strategy for amortizing expensive exploration across deployments: offline distillation compresses multi-task knowledge into transferable initialization that accelerates online convergence without constraining long-term adaptation. Second, our decoupled architecture---separating prompt difficulty learning from model capability profiling---provides a template for building maintainable systems in rapidly evolving model ecosystems where $O(N)$ recalibration becomes prohibitive. Third, our tiered reliability mechanism reconciles conflicting operational objectives (cost minimization vs.\ reliability guarantees) through a single tunable parameter, enabling organizations to dynamically navigate Pareto frontiers without architectural modification.

\paragraph{Limitations and Future Work.}
Several limitations merit future investigation. First, our reliance on a linear reward model (LinUCB) prioritizes $O(1)$ latency and interpretability but may underfit complex, non-linear prompt-model interactions; exploring Neural Contextual Bandits or Kernel methods could improve expressiveness, provided they can maintain millisecond-scale inference. Second, while we mitigate evaluation bias via adversarial baselines (\S\ref{sec:experiments}), our reliance on GPT-4o-as-judge inherently bounds the system's ``ground truth'' to the teacher's worldview; incorporating diverse evaluation protocols (e.g., human preferences, functional correctness tests) would strengthen validity. Third, our current implementation optimizes for single-turn queries; extending to multi-turn dialogues requires modeling conversational state and long-term user satisfaction (e.g., retention) rather than immediate reward.

Future work will explore \emph{Federated Prior Distillation}. Since our ``Shippable Priors'' consist of compact covariance statistics rather than raw data, they offer a unique pathway for multiple organizations to collaboratively improve shared routing intelligence without exposing proprietary query logs. Additionally, we aim to investigate \emph{Adversarial Robustness} to prevent ``Cost-Injection Attacks,'' where malicious actors craft prompts designed to force the router into selecting the most expensive models, thereby draining the deployment budget.

\paragraph{Broader Impact.}
By democratizing access to cost-efficient LLM routing, \ours{} has potential positive societal impact: reducing the economic barrier to deploying AI systems enables broader participation from resource-constrained organizations and researchers. However, the system also introduces risks: aggressive cost optimization might pressure model providers to compete primarily on price, potentially degrading overall ecosystem quality. Additionally, routing decisions encode implicit value judgments about which prompts merit expensive computation; transparency in these allocation mechanisms is essential to prevent systematic biases in who receives high-quality outputs.

\paragraph{Artifact Release.} We release \ours{} as open-source software with pre-trained expert priors at \url{https://github.com/[redacted]/banditgpt}, enabling reproducibility and facilitating adoption in production environments.
